\section*{Capítulo \ref{ch:g7:fractions} --- Frações: comparar e somar}

\begin{solution}{exo:g7:fractions:1}
$\dfrac{2}{5} = \dfrac{8}{20}$; \quad
$\dfrac{18}{24} = \dfrac{3}{4}$; \quad
$\dfrac{7}{3} = \dfrac{28}{12}$; \quad
$\dfrac{5}{9} = \dfrac{20}{36}$.
\end{solution}

\begin{solution}{exo:g7:fractions:2}
$\dfrac{12}{16} = \dfrac34$; \quad
$\dfrac{30}{45} = \dfrac23$; \quad
$\dfrac{27}{36} = \dfrac34$; \quad
$\dfrac{48}{12} = 4$.
\end{solution}

\begin{solution}{exo:g7:fractions:3}
$\dfrac{15}{35} = \dfrac37$ e $\dfrac{21}{49} = \dfrac37$: iguais.

$\dfrac{16}{28} = \dfrac47$ e $\dfrac{20}{36} = \dfrac59$: para
comparar, $\frac47 = \frac{36}{63}$ e $\frac59 = \frac{35}{63}$ --- não
são iguais.
\end{solution}

\begin{solution}{exo:g7:fractions:4}
$\dfrac{7}{11} < \dfrac{9}{11}$ (mesmo \omterm{def:g4:fractions:def}{denominador}).

$\dfrac56 = \dfrac{15}{18} > \dfrac{11}{18}$.

$\dfrac{13}{12} > 1$ enquanto $\dfrac{19}{20} < 1$, logo
$\dfrac{13}{12} > \dfrac{19}{20}$.
\end{solution}

\begin{solution}{exo:g7:fractions:5}
$\dfrac59 + \dfrac29 = \dfrac79$; \quad
$\dfrac{11}{7} - \dfrac47 = \dfrac77 = 1$; \quad
$\dfrac58 + \dfrac78 = \dfrac{12}{8} = \dfrac32$.
\end{solution}

\begin{solution}{exo:g7:fractions:6}
$\dfrac23 + \dfrac56 = \dfrac46 + \dfrac56 = \dfrac96 = \dfrac32$.

$\dfrac{7}{10} - \dfrac25 = \dfrac{7}{10} - \dfrac{4}{10} =
\dfrac{3}{10}$.

$\dfrac34 + \dfrac{5}{12} = \dfrac{9}{12} + \dfrac{5}{12} =
\dfrac{14}{12} = \dfrac76$.
\end{solution}

\begin{solution}{exo:g7:fractions:7}
$3 - \dfrac54 = \dfrac{12}{4} - \dfrac54 = \dfrac74$; \quad
$1 + \dfrac38 = \dfrac{11}{8}$; \quad
$2 - \dfrac76 = \dfrac{12}{6} - \dfrac76 = \dfrac56$.
\end{solution}

\begin{solution}{exo:g7:fractions:8}
$5 \times \dfrac{3}{10} = \dfrac{15}{10} = \dfrac32$; \quad
$\dfrac78 \times 4 = \dfrac{28}{8} = \dfrac72$; \quad
$\dfrac23$ de $45$ é $\dfrac{2 \times 45}{3} = \dfrac{90}{3} = 30$.
\end{solution}

\begin{solution}{exo:g7:fractions:9}
Juntos: $\dfrac14 + \dfrac38 = \dfrac28 + \dfrac38 = \dfrac58$ da
pizza. Sobrou: $1 - \dfrac58 = \dfrac38$.
\end{solution}

\begin{solution}{exo:g7:fractions:10}
Servido: $2 \times \dfrac16 = \dfrac26 = \dfrac13$ L. Resta:
\[
\frac34 - \frac13 = \frac{9}{12} - \frac{4}{12} = \frac{5}{12}
\text{ L}.
\]
\end{solution}

\begin{solution}{exo:g7:fractions:11}
Somar a quantidade positiva $\frac12$ a $\frac35$ tem de dar um
resultado \emph{maior} que $\frac35$. Mas $\frac47 < \frac35$ (de
fato, $\frac47 = \frac{20}{35}$ e $\frac35 = \frac{21}{35}$): a \omterm{def:g1:addition:def}{soma}
alegada seria menor que um dos termos dela --- impossível.
\end{solution}

\begin{solution}{exo:g7:fractions:12}
Passo a passo:
\[
\frac12 + \frac14 = \frac34, \qquad
\frac34 + \frac18 = \frac68 + \frac18 = \frac78, \qquad
\frac78 + \frac1{16} = \frac{14}{16} + \frac1{16} = \frac{15}{16}.
\]
As \omterm{def:g1:addition:def}{somas} parciais $\frac12, \frac34, \frac78, \frac{15}{16}, \dots$
ficam a cada vez na \omterm{def:g3:division:half}{metade} da distância que faltava até $1$: a cada
passo, a parte que falta é cortada ao meio. As \omterm{def:g1:addition:def}{somas} se aproximam de
$1$ sem nunca alcançá-lo.
\end{solution}

\begin{solution}{pb:g7:fractions:1}
\textbf{1.} (a) $\frac12 + \frac14 + \frac14 = \frac24 + \frac14
+ \frac14 = \frac44 = 1$: legal. (b) $\frac14 + \frac18 +
\frac18 + \frac12 = \frac28 + \frac18 + \frac18 + \frac48 =
\frac88 = 1$: legal.

\textbf{2.} $\frac12 + \frac14 + \frac18 = \frac48 + \frac28 +
\frac18 = \frac78$: falta uma colcheia.

\textbf{3.} Um compasso é $1 = \frac{16}{16}$: dezesseis
semicolcheias. Uma colcheia é $\frac18 = \frac{2}{16}$: duas
semicolcheias.

\textbf{4.} Mínima pontuada: $\frac12 + \frac14 = \frac34$. Semínima
pontuada: $\frac14 + \frac18 = \frac28 + \frac18 = \frac38$.

\textbf{5.} A mínima pontuada vale $\frac34$ (questão 4): exatamente um
compasso de valsa. Outros preenchimentos, por exemplo: semínima $+$
semínima $+$ semínima ($\frac14 + \frac14 + \frac14 = \frac34$), ou
semínima $+$ colcheia $+$ colcheia $+$ semínima
($\frac14 + \frac18 + \frac18 + \frac14 =
\frac{2 + 1 + 1 + 2}{8} = \frac68 = \frac34$).

\textbf{6.} $\frac38 + \frac18 + \frac14 = \frac38 + \frac18 +
\frac28 = \frac68 = \frac34$: falta uma semínima ($\frac14$).

\textbf{7.} Colcheia pontuada: $\frac18 + \frac{1}{16} =
\frac{2}{16} + \frac{1}{16} = \frac{3}{16}$. Semínima:
$\frac{4}{16}$. A semínima é mais longa.

\textbf{8.} $\frac12 + \frac18 = \frac48 + \frac18 = \frac58$, ou seja,
$\frac{10}{16}$: dez semicolcheias.

\textbf{9.} $1 - \frac{11}{16} = \frac{16}{16} - \frac{11}{16}
= \frac{5}{16}$ do compasso é silêncio.

\textbf{10.} Os dois exemplos totalizam $\frac{2 + 1 + 1}{16} =
\frac{4}{16}$: legais. Todos os ritmos que enchem quatro semicolcheias
com pedaços de $2$ e de $1$:
\[
2{+}2, \quad 2{+}1{+}1, \quad 1{+}2{+}1, \quad 1{+}1{+}2, \quad
1{+}1{+}1{+}1 :
\]
cinco ritmos.

\textbf{11.} Um ritmo que enche $n$ semicolcheias termina ou com uma
semicolcheia --- e o que vem antes enche $n - 1$ --- ou com uma
colcheia --- e o que vem antes enche $n - 2$. Todo ritmo é contado
exatamente uma vez assim, então
$\text{modos}(n) = \text{modos}(n-1) + \text{modos}(n-2)$. Tabela:
\begin{center}
\renewcommand{\arraystretch}{1.2}
\begin{tabular}{c|cccccccc}
$n$ & $1$ & $2$ & $3$ & $4$ & $5$ & $6$ & $7$ & $8$ \\
\hline
modos & $1$ & $2$ & $3$ & $5$ & $8$ & $13$ & $21$ & $34$ \\
\end{tabular}
\end{center}
Uma mínima ($8$ semicolcheias) pode ser tocada de $34$ maneiras.
($n = 4$ recupera os cinco ritmos da questão 10.)

\textbf{12.} Cada número é a \omterm{def:g1:addition:def}{soma} dos dois anteriores: $3 + 5 = 8$,
$5 + 8 = 13$, $8 + 13 = 21$, $13 + 21 = 34$ --- a regra demonstrada na
questão 11, com mil anos de idade e hoje conhecida como regra de
Fibonacci.

\textbf{13.} $\frac16 + \frac{1}{12} = \frac{2}{12} + \frac{1}{12} =
\frac{3}{12} = \frac14$: a mesma duração de duas colcheias comuns
($\frac18 + \frac18 = \frac14$). A troca do swing é legal.

\textbf{14.} Cada nota da quiáltera dura $\frac{1}{12}$: confira,
$3 \times \frac{1}{12} = \frac{3}{12} = \frac14$
(\cref{prop:g7:fractions:mult}). Para uma mínima: cada nota dura
$\frac16$, já que $3 \times \frac16 = \frac36 = \frac12$.

\textbf{15.} Semínima pontuada $+$ semínima $+$ semínima: em colcheias,
$3 + 2 + 2 = 7$ colcheias $= \frac78$: legal. Mínimas e semínimas valem
$4$ e $2$ colcheias --- os dois números \emph{pares}. Qualquer \omterm{def:g1:addition:def}{soma} de
\omterm{def:g3:numbers:evenodd}{números pares} é par, mas um compasso de sete por oito precisa de $7$
colcheias, um \omterm{def:g3:numbers:evenodd}{número ímpar}: impossível. Para preencher um compasso
\omterm{def:g3:numbers:evenodd}{ímpar}, alguma coisa \omterm{def:g3:numbers:evenodd}{ímpar} --- uma colcheia ou uma nota pontuada ---
tem de aparecer.
\end{solution}
